% TOP_PROBLEM: {"id": 30007118, "title": "Euclidean Steinitz–Bergström bound", "category_id": 6, "category": {"id": 6, "name": "geometry", "display_name": "Geometry", "description": "Euclidean and non-Euclidean geometry, geometric structures.", "slug": "geometry", "order_index": 6, "created_at": "2026-07-31T15:26:25.671Z"}, "status": "open"}
% FIRST_POSTED: 2026-10-06
% ENTRY_KIND: statement_only
% !TeX program = lualatex
\documentclass[11pt]{article}
\usepackage[margin=25mm]{geometry}
\usepackage{fontspec}
\setmainfont{Latin Modern Roman}
\usepackage{amsmath,amssymb,mathtools}
\usepackage{unicode-math}
\setmathfont{Latin Modern Math}
\usepackage{xurl,hyperref}
\hypersetup{colorlinks=true,urlcolor=blue}
\setcounter{secnumdepth}{0}
\setlength{\parindent}{0pt}
\setlength{\parskip}{5pt}
\setlength{\emergencystretch}{3em}
\newcommand{\R}{\mathbb R}
\newcommand{\N}{\mathbb N}
\newcommand{\Z}{\mathbb Z}
\newcommand{\Q}{\mathbb Q}
\newcommand{\C}{\mathbb C}
\newcommand{\D}{\mathbb D}
\newcommand{\F}{\mathbb F}
\newcommand{\sref}[2]{\hyperlink{source:#1:#2}{[S#2]}}
\newcommand{\cataloguescope}[1]{#1}
\begin{document}
\section*{Euclidean Steinitz–Bergström bound}\label{30007118}
\subsection{Definitions and mathematical statement}
\paragraph{Statement recorded in the supplied source.}
Does an absolute C exist such that, for all $d,N\geq 1$ and every indexed family $v_{1},\dots ,v_N$ in the Euclidean unit ball of $R^{d}$ with $\Sigma _i$ $v_i=0$, a permutation $\pi $ satisfies $max_{0\leq k\leq N}\Vert \Sigma _{i\leq k}v_(\pi (i))\Vert _{2}\leq C\surd d$? Repetitions and zeros are permitted, and the same C must be independent of N.
\subsection{Short English statement}
Does an absolute C exist such that, for all $d,N\geq 1$ and every indexed family $v_{1},\dots ,v_N$ in the Euclidean unit ball of $R^{d}$ with $\Sigma _i$ $v_i=0$, a permutation $\pi $ satisfies $max_{0\leq k\leq N}\Vert \Sigma _{i\leq k}v_(\pi (i))\Vert _{2}\leq C\surd d$? Repetitions and zeros are permitted, and the same C must be independent of N.
\subsection{Sources}
\begin{itemize}
\item[S1] Original question source or mathematical context.
\url{http://aimpl.org/hereddiscrep/1/}.
\item[S2] OpenAI, manuscript citation page. Source of the formulation motivating this local problem.
\url{https://github.com/openai/math/blob/adc7f1241b42e322a6451854ab7e4b4c146bf78a/preprints/The-Euclidean-Steinitz-Bergstrom-theorem-September-24-2026/README.md}.
\end{itemize}
\end{document}
